% Generated by support/export-artifacts.py; do not edit.
% Slug: ch42-dependencies
% Source: figures/analysis-dependencies.tex:1
% TikZ source: figures/analysis-dependencies.tex:3
% Original label: \label{fig:ch42-dependencies}
% Original caption (LaTeX source):
%   An arrow runs from a dependency to its reader. A changed callee can invalidate a scalar range proof through the intervening memory facts.
% Caption: An arrow runs from a dependency to its reader. A changed callee can invalidate a scalar range proof through the intervening memory facts.
% Setup: support/compiler-pseudocode.sty (TikZ libraries and shared styles).
\begin{tikzpicture}[node distance=10mm]
\node[basic block] (body) {callee body revision};
\node[basic block,below=of body] (effect) {callee effect summary};
\node[basic block,below=of effect] (load) {caller load facts};
\node[basic block,below=of load] (range) {caller range invariant};
\node[basic block,below=of range] (rewrite) {check-removal decision};
\node[basic block,right=15mm of load] (alias) {alias result};
\draw[cfg edge] (body) -- (effect);
\draw[cfg edge] (effect) -- (load);
\draw[cfg edge] (alias) -- (load);
\draw[cfg edge] (load) -- (range);
\draw[cfg edge] (range) -- (rewrite);
\end{tikzpicture}

% Original figure LaTeX:
% \begin{figure}
% \centering
% \begin{tikzpicture}[node distance=10mm]
% \node[basic block] (body) {callee body revision};
% \node[basic block,below=of body] (effect) {callee effect summary};
% \node[basic block,below=of effect] (load) {caller load facts};
% \node[basic block,below=of load] (range) {caller range invariant};
% \node[basic block,below=of range] (rewrite) {check-removal decision};
% \node[basic block,right=15mm of load] (alias) {alias result};
% \draw[cfg edge] (body) -- (effect);
% \draw[cfg edge] (effect) -- (load);
% \draw[cfg edge] (alias) -- (load);
% \draw[cfg edge] (load) -- (range);
% \draw[cfg edge] (range) -- (rewrite);
% \end{tikzpicture}
% \caption{An arrow runs from a dependency to its reader. A changed callee can invalidate a scalar range proof through the intervening memory facts.}
% \label{fig:ch42-dependencies}
% \end{figure}
